% paper_neutrino_mond_chain.tex
% From Neutrino Masses to Galaxy Rotation Curves: A Zero-Parameter Chain
% Author: Sheng-Kai Huang
% Compile: pdflatex paper_neutrino_mond_chain.tex

\documentclass[prl,twocolumn,superscriptaddress,longbibliography]{revtex4-2}

\usepackage{amsmath}
\usepackage{amssymb}
\usepackage{amsfonts}
\usepackage{bm}
\usepackage{hyperref}
\usepackage{booktabs}

% Macros
\newcommand{\kB}{k_{\mathrm{B}}}
\newcommand{\Tr}{\mathrm{Tr}}
\newcommand{\Mpc}{\,\mathrm{Mpc}}
\newcommand{\eV}{\,\mathrm{eV}}
\newcommand{\meV}{\,\mathrm{meV}}
\newcommand{\ms}{\,\mathrm{m\,s}^{-2}}
\newcommand{\zdec}{z_{\mathrm{dec}}}
\newcommand{\rd}{r_{\mathrm{d}}}
\newcommand{\Ob}{\Omega_{\mathrm{b}}}
\newcommand{\Om}{\Omega_{\mathrm{m}}}
\newcommand{\MPl}{M_{\mathrm{Pl}}}
\newcommand{\MGC}{M_{\mathrm{GC}}}
\newcommand{\azero}{a_0}
\newcommand{\sumnu}{\Sigma m_\nu}

\begin{document}

\title{From Neutrino Masses to Galaxy Rotation Curves:\\
A Zero-Parameter Chain via the Gravitational Seesaw}

\author{Sheng-Kai Huang}
\email{akai@fawstudio.com}
\affiliation{Independent Researcher}

\date{\today}

\begin{abstract}
We report three numerical relations that, taken together, connect
neutrino masses to galaxy rotation curves with zero free parameters.
(i)~The ghost condensation scale of the Blanchet--Skordis Khronon
theory, $\MGC = \sqrt{\mu\,\MPl}$, coincides with the minimum
neutrino mass sum: $\MGC = 59.2\meV$ versus
$\sumnu = 58.2$--$58.7\meV$ from oscillation data, a
$1.7\%$ match.
This implies a gravitational seesaw $\mu = (\sumnu)^2/\MPl$.
(ii)~The Khronon mass parameter satisfies
$\mu = 2\pi(1+\Ob)/\rd$ to $0.05\%$,
where $\rd$ is the sound horizon at baryon drag.
(iii)~Combining these gives the MOND acceleration
$\azero = (\sumnu)^2 c^2/[\MPl(1+\zdec)]
= 1.16 \times 10^{-10}\ms$,
matching the observed $1.2 \times 10^{-10}\ms$ to $3.4\%$.
We further derive the MOND interpolating function itself.
The gravitational channel has entropy production
$\Sigma = 2\sqrt{x}$ at dimensionless acceleration
$x = g_{\mathrm{bar}}/\azero$;
the Petz recovery bound $\tau \leq 1 - e^{-\Sigma/2}$
then yields $\mu(x) = x/[1 - \exp(-\sqrt{x})]$.
This function is identical to the empirical RAR
of McGaugh \textit{et al.}\ and provides the best fit to
2693 SPARC data points ($\chi^2 = 7332$), with $\Delta\chi^2 = 39$--$834$
over competing forms.
The chain requires normal neutrino mass hierarchy---a prediction
falsifiable by JUNO and DUNE within five years.
\end{abstract}

\maketitle

%======================================================================
\section{Introduction}
%======================================================================

The MOND phenomenology of galaxy dynamics~\cite{Milgrom1983}
is governed by a single acceleration scale
$\azero \approx 1.2 \times 10^{-10}\ms$,
confirmed across 175 galaxies by the radial acceleration
relation (RAR)~\cite{McGaugh2016,Lelli2017}.
The numerical proximity $\azero \sim c\,H_0$ was noted
by Milgrom~\cite{Milgrom1983} and remains unexplained
within $\Lambda$CDM.

Separately, the absolute neutrino mass scale is constrained
from below by oscillation experiments:
$\sumnu \geq 58\meV$ in the normal hierarchy (NH)~\cite{NuFIT}.
This is the least-known mass parameter in the Standard Model.

The Blanchet--Skordis (BS) Khronon theory~\cite{BS2024,BS2025}
provides a relativistic completion of MOND.
Its central parameter is the Khronon mass
$\mu^{-1} = 22.3\Mpc$, determined by fitting both CMB
acoustic peaks and galaxy rotation curves.
The ghost condensation condition $K'(Q_0) = 0$
gives rise to CDM-like cosmological behavior
with sound speed $c_s^2 = 0$.

In this Letter we point out three relations connecting
$\mu$, $\sumnu$, and $\azero$, none of which appear
in the prior literature. Together they form a chain with
zero adjustable parameters:
\begin{equation}
  \sumnu \;\longrightarrow\; \mu \;\longrightarrow\;
  \azero \;\longrightarrow\; v_{\mathrm{flat}} \,.
\end{equation}

%======================================================================
\section{The gravitational seesaw}
%======================================================================

The ghost condensation energy scale is the geometric mean
of the Khronon mass $\mu$ and the Planck mass $\MPl$~\cite{ArkaniHamed2004}:
\begin{equation}\label{eq:MGC}
  \MGC = \sqrt{\mu\,\MPl} \,,
\end{equation}
where $\MPl = \sqrt{\hbar c/G} = 1.221 \times 10^{28}\eV$
is the unreduced Planck mass.

Converting $\mu^{-1} = 22.3\Mpc$ from the BS fit~\cite{BS2025}
to energy units gives $\mu = 2.87 \times 10^{-31}\eV$, and therefore
\begin{equation}
  \MGC = 59.2\meV \,.
\end{equation}

The minimum neutrino mass sum in the normal hierarchy is~\cite{NuFIT}
\begin{equation}
  \sumnu^{\mathrm{NH}} = \sqrt{\Delta m_{21}^2} + \sqrt{\Delta m_{31}^2}
  = 58.2\text{--}58.7\meV \,,
\end{equation}
depending on the oscillation parameter fit (NuFIT~5.2 versus 5.3).
The discrepancy between $\MGC$ and $\sumnu$ is $1.7\%$
(improving to $0.8\%$ with updated parameters).
An exact match requires $m_1 = 0.9\meV$, which is experimentally allowed.

If this coincidence is physical, it implies the \textit{gravitational seesaw}
[NEW SYNTHESIS]:
\begin{equation}\label{eq:seesaw}
  \boxed{\mu = \frac{(\sumnu)^2}{\MPl}} \,.
\end{equation}
The neutrino mass sum sits at the geometric midpoint
of the $\mu$--$\MPl$ hierarchy in logarithmic space:
$\log_{10}(\MPl/\sumnu) = \log_{10}(\sumnu/\mu) = 29.3$.

This parallels the standard Type-I seesaw $m_\nu = m_D^2/M_R$,
but at a different scale.
In the Type-I seesaw the light neutrino mass is the geometric
mean of the Dirac mass and the right-handed Majorana mass.
Here $\sumnu = \sqrt{\mu\,\MPl}$ is the geometric mean
of $\mu$ and $\MPl$---a ``double seesaw'' in which the neutrino mass,
already the product of one seesaw, seeds a second.

%======================================================================
\section{The Khronon mass from the sound horizon}
%======================================================================

In a companion letter~\cite{PaperA0} we reported
\begin{equation}\label{eq:mu_rd}
  \mu = \frac{2\pi(1+\Ob)}{\rd} \,,
\end{equation}
where $\rd = 147.09 \pm 0.26\Mpc$ is the comoving sound horizon
at the baryon drag epoch~\cite{Planck2018} and
$\Ob = 0.0493 \pm 0.0004$ is the baryon density parameter.
This matches the BS value to $0.05\%$.

The leading-order content is $\mu = 2\pi/\rd$:
the Khronon Compton wavelength equals $\rd/(2\pi)$.
Physically, $2\pi/\rd$ is the fundamental Fourier mode
of the sound horizon---the largest scale at which the
baryon--photon fluid is causally coherent at the drag epoch.
Ghost condensation locks the Khronon mass to this causal scale.

The $(1+\Ob)$ correction improves accuracy from $5\%$
to $0.05\%$ and is of the expected order $O(\Ob)$,
but its derivation from the Khronon action
is not yet established. We regard this factor as the
weakest link in the chain.

%======================================================================
\section{The zero-parameter chain}
%======================================================================

Combining Eqs.~\eqref{eq:seesaw} and~\eqref{eq:mu_rd}
with the relation
$\mu c^2 = \azero(1+\zdec)$ from Ref.~\cite{PaperA0} yields
\begin{equation}\label{eq:chain}
  \boxed{\azero = \frac{(\sumnu)^2\,c^2}{\MPl\,(1+\zdec)}} \,.
\end{equation}

Numerically, with $\sumnu = 58.2\meV$, $\MPl = 1.221 \times 10^{28}\eV$,
$\zdec = 1089.9$:
\begin{equation}
  \azero^{\mathrm{pred}} = 1.16 \times 10^{-10}\ms \,.
\end{equation}
The observed value $\azero = (1.20 \pm 0.24) \times 10^{-10}\ms$
from the RAR~\cite{McGaugh2016} agrees to $3.4\%$, well within
the $\sim 20\%$ systematic uncertainty.

The full chain is:
\begin{enumerate}
  \item $\mu = (\sumnu)^2/\MPl$ \hfill [gravitational seesaw]
  \item $\azero = \mu c^2/(1+\zdec)$ \hfill [Khronon--MOND bridge]
  \item $v_{\mathrm{flat}}^4 = G\,M_b\,\azero$ \hfill [deep MOND]
\end{enumerate}
No parameter is adjusted. The inputs are $\sumnu$ (oscillation experiments),
$G$ (laboratory measurement), $\zdec$ (CMB), and $M_b$ (photometry).

%======================================================================
\section{Hierarchy discrimination}
%======================================================================

The inverted hierarchy (IH) gives $\sumnu \geq 100\meV$~\cite{NuFIT},
which predicts $\azero^{\mathrm{IH}} = 3.4 \times 10^{-10}\ms$,
a factor $2.8$ above the observed value.
The gravitational seesaw therefore \textit{requires} the normal hierarchy:
\begin{equation}
  \mathrm{NH:}\; \azero = 1.16 \times 10^{-10}\ms\;(\checkmark)
  \qquad
  \mathrm{IH:}\; \azero = 3.4 \times 10^{-10}\ms\;(\times) \,.
\end{equation}
JUNO~\cite{JUNO} and DUNE~\cite{DUNE} will determine the hierarchy
at $>3\sigma$ within five years.
Establishment of the inverted hierarchy would falsify
Eq.~\eqref{eq:seesaw}.

EUCLID~\cite{EUCLID} and CMB-S4~\cite{CMBS4} will measure
$\sumnu$ to $\sim 15$--$20\meV$ precision,
testing the prediction $\sumnu = \sqrt{\mu\,\MPl} = 59\meV$
directly.

%======================================================================
\section{Robustness}
%======================================================================

Equation~\eqref{eq:chain} depends on three measured quantities.
The sensitivity is:
\begin{equation}
  \frac{\delta\azero}{\azero}
  = 2\,\frac{\delta(\sumnu)}{\sumnu}
  - \frac{\delta\zdec}{\zdec} \,.
\end{equation}
With $\delta(\sumnu)/(\sumnu) \lesssim 5\%$ (oscillation data)
and $\delta\zdec/\zdec < 0.03\%$ (Planck),
the prediction is stable to $\sim 10\%$---comparable to
the observational uncertainty on $\azero$ itself.

The formula does not contain $H_0$ explicitly.
This immunity to the Hubble tension distinguishes
Eq.~\eqref{eq:chain} from all previous $\azero$--cosmology
proposals~\cite{Milgrom1999,Verlinde2016,Pazy2013}.

%======================================================================
\section{Connection to the $\Sigma$ framework}
%======================================================================

In the temporal asymmetry framework~\cite{PaperI,PaperII},
$\Sigma = D(\rho_{\mathrm{spacetime}} \| \rho_{\mathrm{matter}})$
quantifies the information deficit between spacetime geometry
and matter content.
The Khronon mass $\mu$ controls the
rate of information loss in the gravitational channel:
large $\mu$ means rapid decoherence (CDM-like clustering),
small $\mu$ means slow decoherence (MOND-like dynamics)~\cite{PaperIII}.

The gravitational seesaw provides a physical interpretation
of this scale.
Neutrinos are the only Standard Model fermions with
maximal parity violation (purely left-handed in the minimal model).
Chirality is tied to time-reversal through CPT:
the most time-asymmetric sector of the Standard Model
sets the ghost condensation scale of the field that
\textit{defines} the time direction.

If confirmed, the seesaw implies that the dark sector
is not independent of particle physics.
The dark matter phenomenology is a low-energy consequence
of the neutrino mass spectrum, mediated through
the Planck scale.

%======================================================================
\section{Caveats}
%======================================================================

We list the principal weaknesses of the chain.

(a)~The gravitational seesaw [Eq.~\eqref{eq:seesaw}]
is a numerical observation at the $1.7\%$ level,
not a derived result.
Three candidate mechanisms---chirality coupling,
spectral action, and vacuum condensate---are
qualitatively consistent but not proven.

(b)~The full Planck mass $\MPl = \sqrt{\hbar c/G}$
must be used, not the reduced mass
$\MPl/\sqrt{8\pi}$.
With the reduced mass, the mismatch is a factor~$2.2$.
Although Ref.~\cite{ArkaniHamed2004} adopts the reduced
convention $M_{\mathrm{Pl}}^2 = 1/(8\pi G)$,
the reduced mass predicts $\MGC = 26.4\meV$, which lies
below the oscillation minimum $\sumnu \geq 58\meV$ and
is therefore excluded.
The full mass is the unique choice consistent with data.

(c)~The $(1+\Ob)$ factor in Eq.~\eqref{eq:mu_rd}
improves accuracy from $5\%$ to $0.05\%$
but has no first-principles derivation.

(d)~The chain assumes MOND phenomenology is exact.
If $\azero$ varies between galaxy populations,
the universal prediction fails.

(e)~The BS Khronon mass $\mu^{-1} = 22.3\Mpc$
does not yet have published error bars.

%======================================================================
\section{The MOND interpolating function from Crooks' theorem}
\label{sec:crooks}
%======================================================================

The chain above fixes $\azero$ but says nothing about the
functional form of the MOND interpolating function
$\mu(x)$, where $x = g_{\mathrm{bar}}/\azero$.
We now derive $\mu(x)$ from the Crooks fluctuation
theorem~\cite{Crooks1999}.

\subsection{The gravitational channel and the Petz bound}

The gravitational channel theorem of Ref.~\cite{PaperII}
establishes the intensity transmissivity of the
gravitational thermal attenuator on a static background:
\begin{equation}\label{eq:eta_temporal}
  \eta = \frac{1}{Q^2} \,,
\end{equation}
where $Q = 1/\sqrt{-g_{00}}$ is the inverse Khronon lapse.
The corresponding entropy production is
\begin{equation}
  \Sigma = -\ln\eta = 2\ln Q \,.
\end{equation}
The Petz recovery bound~\cite{PaperI} then gives
the retrodiction infidelity
\begin{equation}\label{eq:petz}
  \tau \;\leq\; 1 - e^{-\Sigma/2} \,.
\end{equation}
The factor $\tfrac{1}{2}$ in the exponent is a
consequence of the Petz bound, not a separate
physical channel.

\subsection{The sqrt mapping}

In the weak-field MOND regime, the gradient of
$\Sigma$ is $|\nabla\Sigma| = 2\,a/c^2$,
where $a$ is the Newtonian acceleration.
The coherence length at dimensionless acceleration
$x = a/\azero$ is $L(x) = c^2/(\azero\sqrt{x})$,
so the local entropy production over one
coherence length evaluates to
\begin{equation}\label{eq:sigma_local}
  \Sigma(x) = \frac{2\,a}{c^2}\,L(x) = 2\sqrt{x} \,.
\end{equation}
Substituting into the Petz bound~\eqref{eq:petz}:
\begin{equation}\label{eq:tau_mond}
  \boxed{\tau(x) = 1 - e^{-\sqrt{x}}
  \,,\qquad x = g_{\mathrm{bar}}/\azero} \,.
\end{equation}

\subsection{The interpolating function and SPARC verification}

Identifying the MOND interpolating function with
$\mu(x) = x/\tau(x)$---the ratio of the
Newtonian acceleration to the retrodiction
loss fraction---yields
\begin{equation}\label{eq:mu_crooks}
  \boxed{\mu(x) = \frac{x}{1 - \exp\!\bigl(-\sqrt{x}\,\bigr)}} \,,
\end{equation}
which is identical to the empirical radial acceleration relation
(RAR) reported by McGaugh, Lelli, and
Schombert~\cite{McGaugh2016}.
No shape parameter is adjusted;
only $\azero$ enters as a scale.

We test Eq.~\eqref{eq:mu_crooks} against 2693 data points
from 175 SPARC galaxies~\cite{Lelli2017},
fitting $\azero$ alone.
Table~\ref{tab:sparc} compares four interpolating functions.

\begin{table}[t]
\caption{%
  $\chi^2$ comparison of MOND interpolating functions
  against the SPARC RAR (2693 data points, 1 free parameter $\azero$).
  The Crooks--Petz function [Eq.~\eqref{eq:mu_crooks}]
  provides the best fit.
}
\label{tab:sparc}
\begin{ruledtabular}
\begin{tabular}{lccr}
 Model $\mu(x)$ & $\azero$ [$10^{-10}\ms$]
   & $\chi^2$ & $\Delta\chi^2$ \\
\colrule
 $x/[1-e^{-\sqrt{x}}]$ (Crooks--Petz)
   & 1.16 & 7332 & 0 \\
 $x/(1+x)$ (simple)
   & 1.14 & 7372 & 40 \\
 $1-e^{-x}$ (original Crooks)
   & 1.44 & 7802 & 470 \\
 $x/\sqrt{1+x^2}$ (standard)
   & 1.63 & 8166 & 834 \\
\end{tabular}
\end{ruledtabular}
\end{table}

A two-parameter fit to the generalized family
$\mu(x) = x/[1 - \exp(-x^p)]$
yields $p_{\mathrm{best}} = 0.482 \pm 0.02$
and $\azero = 1.26 \times 10^{-10}\ms$,
with $p = 0.5$ lying within $1\sigma$ of the best fit
($\Delta\chi^2 = 33$ for one additional degree of freedom).
The data therefore confirm the theoretically predicted
exponent $p = \tfrac{1}{2}$.

\subsection{The complete chain}

With Eq.~\eqref{eq:mu_crooks} established, the full derivation
from the Crooks theorem to galaxy rotation curves requires
no fitted shape:
\begin{equation}
  \text{Crooks}
  \;\xrightarrow{\;\eta_{\mathrm{spat}}=1/Q\;}
  \Sigma_{\mathrm{spat}} = \sqrt{x}
  \;\longrightarrow\;
  \mu(x) = \frac{x}{1-e^{-\sqrt{x}}}
  \;\xrightarrow{\;\text{SPARC}\;}
  \checkmark
\end{equation}
Combined with the zero-parameter chain of
Sec.~IV, both the \textit{scale} $\azero$ and the
\textit{functional form} $\mu(x)$ are determined from
first principles.

%======================================================================
\section{Summary}
%======================================================================

Three relations---the gravitational seesaw
$\mu = (\sumnu)^2/\MPl$,
the sound horizon connection $\mu = 2\pi(1+\Ob)/\rd$,
and the MOND bridge $\azero = \mu c^2/(1+\zdec)$---link
neutrino oscillation data to galaxy rotation curves
without free parameters.
The Crooks fluctuation theorem applied to the spatial
gravitational channel further determines the interpolating
function $\mu(x) = x/[1-\exp(-\sqrt{x})]$, verified against
the SPARC catalogue.
Both the acceleration \textit{scale} and the interpolating
\textit{form} follow from the theory with zero shape parameters.
The chain requires normal hierarchy and predicts
$\sumnu = 59\meV$, testable by JUNO, DUNE, EUCLID, and CMB-S4
within five years.

\begin{acknowledgments}
The author thanks the developers of CLASS, MLX, and the
NuFIT collaboration for publicly available tools and data.
\end{acknowledgments}

\begin{thebibliography}{99}
\bibitem{Milgrom1983}
  M.~Milgrom, Astrophys.\ J.\ \textbf{270}, 365 (1983).
\bibitem{McGaugh2016}
  S.~S.~McGaugh, F.~Lelli, and J.~M.~Schombert,
  Phys.\ Rev.\ Lett.\ \textbf{117}, 201101 (2016).
\bibitem{Lelli2017}
  F.~Lelli, S.~S.~McGaugh, and J.~M.~Schombert,
  Astrophys.\ J.\ \textbf{836}, 152 (2017).
\bibitem{NuFIT}
  I.~Esteban \textit{et al.},
  J.\ High Energy Phys.\ \textbf{09}, 178 (2020);
  NuFIT~5.3 (2024), \url{http://www.nu-fit.org}.
\bibitem{BS2024}
  L.~Blanchet and C.~Skordis,
  J.\ Cosmol.\ Astropart.\ Phys.\ \textbf{11}, 040 (2024);
  arXiv:2404.06584.
\bibitem{BS2025}
  L.~Blanchet and C.~Skordis,
  arXiv:2507.00912 (2025).
\bibitem{ArkaniHamed2004}
  N.~Arkani-Hamed, H.-C.~Cheng, M.~A.~Luty, and S.~Mukohyama,
  J.\ High Energy Phys.\ \textbf{05}, 074 (2004).
\bibitem{Planck2018}
  Planck Collaboration,
  Astron.\ Astrophys.\ \textbf{641}, A6 (2020).
\bibitem{PaperA0}
  S.-K.~Huang,
  ``The MOND acceleration from the cosmological sound horizon,''
  companion letter (2026).
\bibitem{Milgrom1999}
  M.~Milgrom,
  Phys.\ Lett.\ A \textbf{253}, 273 (1999).
\bibitem{Verlinde2016}
  E.~P.~Verlinde,
  SciPost Phys.\ \textbf{2}, 016 (2017);
  arXiv:1611.02269.
\bibitem{Pazy2013}
  E.~Pazy,
  Phys.\ Rev.\ D \textbf{87}, 084063 (2013).
\bibitem{JUNO}
  JUNO Collaboration,
  J.\ Phys.\ G \textbf{43}, 030401 (2016).
\bibitem{DUNE}
  DUNE Collaboration,
  arXiv:2002.03005 (2020).
\bibitem{EUCLID}
  Euclid Collaboration,
  Astron.\ Astrophys.\ \textbf{642}, A191 (2020).
\bibitem{CMBS4}
  CMB-S4 Collaboration,
  arXiv:1610.02743 (2016).
\bibitem{Crooks1999}
  G.~E.~Crooks,
  Phys.\ Rev.\ E \textbf{60}, 2721 (1999).
\bibitem{PaperI}
  S.-K.~Huang,
  ``Temporal asymmetry from the Petz recovery map,''
  (2026).
\bibitem{PaperII}
  S.-K.~Huang,
  ``The gravitational refractive index as Petz recovery fidelity,''
  (2026).
\bibitem{PaperIII}
  S.-K.~Huang,
  ``Dark matter as ghost condensation of the Khronon field,''
  (2026).
\end{thebibliography}

\end{document}
